\section{\texorpdfstring{T1: The Expressiveness Boundary\\(Oracle Enforceability)}{T1: The Expressiveness Boundary (Oracle Enforceability)}}\label{sec:t1}

Call $g$ a \emph{$(\Th;\Pi)$-gate} if it runs a deterministic keyed-counter register automaton
(Def.~\ref{def:safety}) with $\Th$-guards and queries to the fixed predicate bits $\Pi$ once, online
and in commit order, over the committed \emph{effect} history, and decides each proposal by the one
transition it takes from the configuration reached, labeled $\allow$, $\block$, or $\subst(t)$ for a
term $t$ over the proposal's fields, the registers, and constants;
Remark~\ref{rem:general-gates} treats general computable gates.

\begin{theorem}[Oracle characterization; classical \citep{ligatti2005edit,ligatti2009nonsafety}, instantiated]\label{thm:t1}
Fix any oracle interface $\Pi$ (Def.~\ref{def:judge}; with action-local $\Pi$ the bits cannot replace the automaton's memory, Proposition~\ref{prop:strict}(ii)) and a nonempty policy $P$. Then $P$ is precisely enforceable
by a deterministic $(\Th;\Pi)$-gate if and only if it is safety and
$\Good(P)\in\RA[\Th;\Pi]$. Enforcement is witnessed by simulating that recognizer
and refusing exactly the one-step extensions it rejects; each step evaluates one
enabled transition in the given guard theory.
\end{theorem}

\noindent\emph{Proof idea.} ($\Leftarrow$) Simulate the recognizer on the committed history and
block exactly the extensions it rejects. ($\Rightarrow$) If $P$ is not safety, some violating
$x$ is a limit point of $P$; transparency on its compliant extensions forces the deterministic
gate to allow every step of $x$, contradicting soundness. Simulating the gate on finite words
yields a recognizer of $\Good(P)$ (Supplement~A.2).

\paragraph{Running example}
P3 needs one control flag set by $\mathsf{revoke}$, and P2 one monotone counter per
key $(k,w)$ with guard $\mathit{spent}_{k,w}+v\le c$, so a gate enforces both
without semantic predicates and never blocks a compliant send or transfer. P1's
guard, ``the body contains PII,'' is semantic: a gate enforces P1 only if that
predicate is in $\Pi$; with a fallible judge in its place, T2 governs.

\begin{remark}[General gates]\label{rem:general-gates}
An arbitrary computable gate with the same oracle interface enforces a nonempty $P$ if and only
if $P$ is safety with good-prefix membership decidable relative to $\Pi$: the
truncation-automaton characterization of \citet{ligatti2005edit,ligatti2009nonsafety} (already over countably infinite action sets), with blocking in place of
halting and with oracle-relative decidability. $(\Th;\Pi)$-gates
are its register-recognizable fragment.
\end{remark}

\begin{proposition}[A strict hierarchy]\label{prop:strict}
Give every enforcer the same oracle bits $\Pi$.
(i)~Gates, which enforce precisely, enforce strictly fewer policies than edit automata under effective
enforcement \citep{ligatti2005edit} (an $\omega$-policy counts as edit-enforceable when some extension of it to
finite executions is, Supplement~A.2); the gap is exactly buffering, which irreversibility removes. Witness: ``every $\mathsf{pay}$ is eventually followed by a
$\mathsf{confirm}$'' is a renewal property an edit automaton enforces by buffering
$\mathsf{pay}$, but transparency on the compliant streams $\mathsf{pay}^{k}\mathsf{confirm}^{\omega}$
forces a gate to allow every irreversible $\mathsf{pay}$, so on $\mathsf{pay}^{\omega}$ it
commits a violation no later decision can retract.
(ii)~If $\Pi$ is action-local and recipients range over an infinite equality-only data sort, then no
$(\Th;\Pi)$-gate without key maps (whatever its registers, global counters, and counter guards and
updates) enforces the per-recipient cap of one, which a $(\Th_0;\Pi)$-gate enforces with one key-local
flag. Witness: after $r{+}1$ allowed sends to distinct recipients that avoid the gate's constants and
share oracle bits, some recipient $d$ is in none of the gate's $r$ registers, and a fresh such recipient
$e$ satisfies exactly the same guards, although a second send to $d$ is a violation and a send to $e$ is
not.
\end{proposition}

Part~(i) is the classical truncation/edit separation \citep{ligatti2005edit,ligatti2009nonsafety},
restated to locate gates, and holds for every $\Pi$; part~(ii) is where the restriction on $\Pi$ matters,
since with the history bit $\mathbf{1}[\tau\cdot a\in\Good(P)]$ a register-free gate enforces the cap
(\S\ref{sec:model}). The witness of part~(i) is a common design: a deferred-commit payment protocol
needs buffering, so the deferral must live in the tool interface, not in the gate. Over finite
alphabets, enforcers for reactive systems must release an output at every reaction, so not every safety property
is enforceable there \citep{pinisetty2017synchronous,anand2025prompt}; our gate can always commit
the no-op. Over data, determinism is a genuine restriction \citep{aceto2025monitorability}, and
relational policies are hyperproperties outside this class
\citep{clarkson2010hyperproperties,pandey2026local}.

Substitution ($\subst(a')$) leaves the class unchanged, and so does escalation if the escalation
oracle is part of $\Pi$: the limit-point argument and the never-alter recognizer apply verbatim. With an oracle, a transparent gate never alters a compliant proposal,
so substitution cannot improve $U$; with a judge, committing a \emph{prefix-safe} rewrite ($ua'$
good whenever $u$ is) instead of blocking lowers the graded false-block loss to at most
$\ell_{\rm sub}\alpha(\theta)$ at unchanged miss, where $\ell_{\rm sub}\in[0,1]$ is a rewrite's loss relative to a block's unit loss. Escalation fits only as an oracle that returns
prefix-safe actions (Supplement~A.2).
